\documentclass[10pt,a4paper]{amsart}
\usepackage[margin=19mm]{geometry}
\usepackage{amsmath,amssymb,mathtools}
\usepackage[scr=rsfs,cal=euler]{mathalfa}
\usepackage{xfrac}
\usepackage[unicode,hidelinks,bookmarks=false]{hyperref}
\newcommand{\ie}{i.e.\ }
\newcommand{\cf}{cf.\ }
\newcommand{\resp}{respectively\ }
\newcommand{\Subsection}{\subsection}
\providecommand{\nobreakdash}{\nobreak\hbox{-}}
\setlength{\emergencystretch}{2em}
\multlinegap0pt
\newcommand\QQ{\mathbb{Q}}
\usepackage{tikz}
\usetikzlibrary{cd}
\usepackage{mathtools}
\usepackage{amssymb}
\usepackage{xcolor}
\usepackage{enumerate}
\usepackage{bm}
\usepackage{cancel}
\usepackage[normalem]{ulem}
\usepackage[shortlabels]{enumitem}
\usepackage{tikz}
\newtheorem{lemma}{Lemma}
\newcommand\ZZ{\mathbb{Z}}
\title{$\QQ_p$-homotopy types and applications to topology and algebraic geometry}
\author{Runjie Hu, Guozhen Wang}
\date{Source: arXiv:2608.25389v1 (CC BY 4.0)}
\begin{document}
\maketitle

\noindent\emph{Source and licence.} $\QQ_p$-homotopy types and applications to topology and algebraic geometry. Version arXiv:2608.25389v1 (CC BY 4.0), distributed by arXiv under the Creative Commons Attribution 4.0 International licence, http://creativecommons.org/licenses/by/4.0/, as recorded in the arXiv licence metadata for this exact version and shown on the abstract page https://arxiv.org/abs/2608.25389v1. Full licence notice: \texttt{licenses/arxiv-2608.25389-CC-BY-4.0.txt}.\par\smallskip

\noindent\textit{Editorial scope.} Counted unit: the article's lemma lifting (source0.tex lines 899--901). The article's complete original proof of it is delivered with it (source0.tex lines 903--920, one \texttt{proof} environment of 18 lines). No mathematics is written, repaired or re-derived: the statement and the proof are the article's own lines, in the article's own notation, and no retained line is substituted or re-flowed.  No further source-local context is needed: every cross-reference of every retained line resolves inside the two delivered spans.  Nothing was omitted from any retained span, so the package carries no omission record.  No retained line contains a citation, so the wrapper carries no bibliography and no bibliography entry.  The retained lines contain the article's own diagram code, which the wrapper typesets with the standard tikz packages; no image file is delivered and the package declares no asset dependency.  Editorial formatting only, in addition: an AMS \texttt{amsart} preamble that restates the article's own declarations and macros used by the retained lines, namely the environments \texttt{lemma} on separate counters, together with the standard packages those lines need.  The retained environments print in this document as Lemma 1 in order of appearance, whereas the article prints the same items with the numbering of its own full document.  The complete native source of the article as distributed in the arXiv e-print for this exact version is bundled unchanged as \texttt{sources/208576/source0.tex}.
\begin{lemma}\label{Lem: local spaces have an integral lifting}
Let $T$ be a set of prime numbers. Every nilpotent $T$-local finite type CW complex is the $T$-localization of a nilpotent finite type CW complex.  
\end{lemma} 

\begin{proof}
We argue by induction on a refined principal Postnikov tower. It suffices to prove the following: for a nilpotent finite type CW complex $X$ and a principal fibration $K(G_{(T)},n)\xrightarrow{i} Y_{(T)}\xrightarrow{q} X_{(T)}$, where $G$ is a finitely generated abelian group, $G_{(T)}=G\otimes_{\ZZ}\ZZ_{(T)}$, and $X_{(T)}$  is a $T$-localization of $X$, there exists a nilpotent finite type CW complex $Y$, together with a map of principal fibrations as in the following diagram, such that the vertical arrows are $T$-localizations.
\[
\begin{tikzcd}
    K(G,n) \arrow[r] \arrow[d] & Y \arrow[r] \arrow[d] & X \arrow[d] \\
    K(G_{(T)},n) \arrow[r,"i"] & Y_{(T)} \arrow[r,"q"] & X_{(T)}
\end{tikzcd}
\]
Equivalently, it suffices to construct the following diagram, where $b$ is induced by the principal fibration $K(G_{(T)},n)\xrightarrow{i} Y_{(T)}\xrightarrow{q} X_{(T)}$, the upper vertical arrows are $T$-localizations, and the lower vertical arrows are homotopy equivalences.
\[
\begin{tikzcd}
    Y \arrow[r] \arrow[d] & X \arrow[r] \arrow[d] & K(G,n+1) \arrow[d] \\
     Y'_{(T)} \arrow[r,"q'"] \arrow[d,"\simeq"] & X'_{(T)} \arrow[r,"b'"] \arrow[d,"\simeq"] & K(G_{(T)},n+1) \arrow[d,"\simeq","\alpha"'] \\
     Y_{(T)} \arrow[r,"q"] & X_{(T)} \arrow[r,"b"]  & K(G_{(T)},n+1) 
\end{tikzcd}
\]
The natural isomorphism $H^{n+1}(X_{(T)};G_{(T)})\cong H^{n+1}(X;G)\otimes_{\ZZ} \ZZ_{(T)}$ implies that there exist a class $[b']\in H^{n+1}(X;G)$ and an integer $k$ not divisible by any prime in $T$ such that $k\cdot [b']=[b]\in H^{n+1}(X_{(T)};G_{(T)})$. Let $\alpha$ be the self-map of $K(G_{(T)},n+1)$ induced by multiplication by $\frac{1}{k}$ on $ G_{(T)}$. Define the middle row by pulling back the bottom fibration along $\alpha$. The localization map $X\rightarrow X_{(T)}$ induces a map $X\rightarrow X'_{(T)}$ which makes the upper-right square commute up to homotopy. Since $X'_{(T)}\rightarrow X_{(T)}$ is a homotopy equivalence, $X\rightarrow X'_{(T)}$ is also a $T$-localization. Define $Y$ to be the homotopy fiber of $b':X\rightarrow K(G,n+1)$. The resulting map $Y\rightarrow Y'_{(T)}$ has the required properties, completing the proof.
\end{proof}

\end{document}
